% ============ 34 估计的几何 (原书 pp.197-199) ============
\chapter{估计的几何}\label{sec:34}
\index{指数和}
\index{Kloosterman 和}
\index{Gauss 和}
\index{单值}
\index{Deligne 定理}

Weil 猜想表明：有限域上方程组解的数量由方程提升到特征零后定义的复簇的拓扑性质控制.

本节解释指数(及类似)和的大小估计如何反映某些层上的单值作用.

例如，设 $X$ 是 $\mathbb{F}_q$ 上的光滑射影曲面，$\varphi \colon X \to \mathbb{P}^{1}$ 是正则映射，其纤维 $X_t$ 除有限多个 $t \in \mathbb{F}$ 外都是椭圆曲线.对 $t \in \mathbb{P}^{1}(\mathbb{F}_q)$，设
\[
\#X_t(\mathbb{F}_q) = q - e(t) + 1.
\]
若忽略奇异纤维，则 Weil 猜想告诉我们
\[
|2e(t)| < 2\sqrt{q}.
\]
对
\[
\sum_{t} e(t), \qquad \sum_{t} e(t)^{2}, \qquad \sum_{t} e(t)e(t+u), \qquad \sum_{t_1 + t_2 = t_3 + t_4} e(t_1)e(t_2)e(t_3)e(t_4)
\]
是否有类似的估计？求和遍历 $t \in \mathbb{F}_q$.解析数论学家有猜测这种和的估计的启发式方法.例如，在 $\sum e(t)$ 中有 $q$ 项，每项的大小约为 $2\sqrt{q}$，因此平凡地该和 $\leq Cq^{\frac{3}{2}}$($C$ 是某个与 $p$ 的幂 $q$ 无关的常数).然而，除非族常值，否则预期 $e(t)$ 随机变化，这提示
\[
\left| \sum_{t} e(t) \right| \leq Cq.
\]
这个特定的不等式有初等证明，但一般地，想要的结果无法用解析数论的方法得到.Deligne 的定理给出了得到这种估计的非常有力的方法.把和解释为曲线上层的上同调群上的 Frobenius 算子的迹，并作为层的几何(单值)与 Deligne 定理理解的推论得到估计.例如，可以这样得到如下结果.

\begin{theorem}\label{thm:34.1}
设 $e(t)$ 如上.
\begin{enumerate}
  \item 若族的 $j$ 不变量不常值，则
  \[
  \left| \sum_{t} e(t) \right| \leq (\beta_2(X) - 2)q,
  \]
  其中 $\beta_2(X)$ 是 $X$ 的第二个 Betti 数.
  \item 若 $j$ 不常值，则
  \[
  \sum_{t} e(t)^{2} = q^{2} + O(q^{\frac{3}{2}}).
  \]
  \item 设 $S$ 是使 $X_t$ 奇异的 $t \in \mathbb{F}$ 的集合；若集合 $S$ 与 $\{s - u \mid s \in S\}$ 不相交，则
  \[
  \sum_{t} e(t)e(t+u) = O(q^{\frac{3}{2}}).
  \]
\end{enumerate}
\end{theorem}

第四个和困难得多，但下面已知.假定 $\varphi$ 使得 $X_t$ 是射影曲面
\[
X^{3} + Y^{3} + Z^{3} = nW^{3}
\]
与平面
\[
\alpha X + \beta Y + \gamma Z = tW
\]
的交.则，在 $q$ 的某些假设下，
\[
\sum_{t_1 + t_2 = t_3 + t_4} e(t_1)e(t_2)e(t_3)e(t_4) \geq 2q^{4} - Bq^{7/2}
\]
对某个\footnote{该不等式应 C. Hooley 的要求证明，并使他得到了整数表示为 3 个立方与 2 个平方之和的方式数的渐近估计，这是他不成功地工作了 20 多年的问题.见他在 International Congress of Mathematicians, Warsaw 1983 的大会报告，该报告没有承认平展上同调在他的最终成功中起的关键作用.}常数 $B$(Milne, J., \emph{Estimates from \'etale cohomology}, Crelle 328, 1981, 208--220).

关于平展上同调对各种和的估计的应用的全面论述，见：
\begin{quote}
Katz, N., \emph{Sommes exponentielles}. Course taught at the University of Paris, Orsay, Fall 1979. With a preface by Luc Illusie. Notes written by G\'erard Laumon. Ast\'erisque, 79. Soci\'et\'e Math\'ematique de France, Paris, 1980. 209 pp.

Katz, N., \emph{Gauss sums, Kloosterman sums, and monodromy groups}. Annals of Mathematics Studies, 116. Princeton University Press, Princeton, NJ, 1988. x+246 pp.
\end{quote}

数论估计与层的单值之间的对应可以双向使用：N. Katz 用 Davenport 与 Lewis 关于有限域上多项式解的一些估计证明了如下定理(\emph{Monodromy of families of curves: applications of some results of Davenport-Lewis}. Seminar on Number Theory, Paris 1979--80, pp. 171--195, Progr. Math., 12, Birkh\"auser, Boston, Mass., 1981):

\begin{theorem}\label{thm:34.2}
设 $f(X, Y) \in \mathbb{C}[X, Y]$ 是两个变量的多项式.假定对未定元 $a, b, c$，仿射曲线
\[
f(X, Y) + aX + bY + c = 0
\]
在域 $\mathbb{C}(a, b, c)$ 上的完备非奇异模型有亏格 $g \geq 1$.则对 $\mathbb{A}^{3}_{\mathbb{C}}$ 的任何非空 Zariski 开集 $S$(在其上完备非奇异模型``漂亮地''延拓为态射 $\bar{f} \colon \mathcal{C} \to S$)，$S$ 的基本群在 $R^{1}\bar{f}_{*}\mathbb{Q}$ 的一般茎上绝对不可约作用($R^{1}\bar{f}_{*}\mathbb{Q}$ 是对复拓扑的高阶直像).
\end{theorem}

数论估计与平展层的几何之间的相互作用引人入胜，但要充分欣赏它，需要同时理解解析数论与平展上同调.

\subsection*{Katz 1980 书评(MR 82m:10059)}

这些讲义用平展上同调证明两个关于指数和大小 的非常一般性质的定理.

对任何 $q = p^{n}$，设 $\psi \colon \mathbb{F}_q \to \mathbb{C}^{\times}$ 是加性特征 $a \mapsto \exp\left( \frac{2\pi i}{p} \operatorname{Tr}_{\mathbb{F}_q/\mathbb{F}_p} a \right)$.设 $V$ 是 $\mathbb{F}_p$ 上的簇(或有限型概形)，$f \colon V \to \mathbb{A}^{1}_{\mathbb{F}_p}$ 是 $V$ 上的正则函数.与 $V$，$f$ 与 $q = p^{n}$ 相联系的指数和是 $S_q(V, f) = \sum \psi(f(x))$，其中求和遍历所有 $x \in V(\mathbb{F}_q)$.例如，若 $V$ 由 $XY = a$ 定义，$f$ 是函数 $(x, y) \mapsto x + y$，则
\[
S_p(V, f) = \sum_{xy = a} \exp\left( \frac{2\pi i}{p}(x + y) \right) = \sum_{x \neq 0} \exp\left( \frac{2\pi i}{p}\left( x + \frac{a}{x} \right) \right)
\]
是 Kloosterman 和.

第一个定理处理如下问题：假定 $V$ 由系数在 $\mathbb{Z}$ 中的方程定义，$f$ 是映射 $V \to \mathbb{A}^{1}_{\mathbb{Z}}$，因此对所有 $p$，有由 $(V, f)$ 模 $p$ 约化得到的 $\mathbb{F}_p$ 上的对 $(V_p, f_p)$；那么能否一致地界定和 $S_q(V_p, f_p)$？给出的答案是：对固定的 $(V, f)$，存在常数 $A$，使得对所有足够大的 $p$ 与所有被这样的 $p$ 整除的 $q$，$|S_q(V_p, f_p)| \leq Aq^{N}$，其中 $N$ 是 $f$ 的几何纤维的最大维数，且假定一般纤维几何不可约或维数 $< N$.证明用标准论证把 $S_q$ 解释为某个平展上同调空间上 Frobenius 自同态的迹的交替和，然后用平展上同调的一般结果与 P. Deligne [Inst. Hautes \'Etudes Sci. Publ. Math. No. 52 (1980), 137--252] 的结果分别界定空间的维数与权.

第二个定理在几何特别容易处理的情形为固定的 $p$ 提供显式常数 $A$.例如考虑由系数在 $\mathbb{F}_p$ 中的次数为 $D$ 的齐次方程 $P(X_0, X_1, X_2, X_3) = 0$ 定义的非奇异射影曲面 $V_0$.设 $H$ 是由 $X_0 = 0$ 定义的超平面，$F$ 是线性形式 $a_1 X_1 + a_2 X_2 + a_3 X_3$.假定 $(D, p) = 1$，$C = V_0 \cap H$ 是次数 $D$ 的光滑曲线，且平面 $F = 0$ 横截相交 $C$.设 $V$ 是仿射曲面 $V_0 \setminus C$(由 $P(1, X_1, X_2, X_3) = 0$ 定义)，$f$ 是函数 $(x_1, x_2, x_3) \mapsto \sum a_i x_i \colon V \to \mathbb{A}^{1}_{\mathbb{F}_p}$.则由 Deligne(同上)与 Grothendieck-Ogg-Shafarevich 公式相当直接地推出 $|S_q(V, f)| \leq D(D-1)^{2}q$.在讲义中，这个结果推广到 $V_0$ 是任何光滑射影簇，$H$ 是由方程 $s = 0$ 定义的超平面，$f$ 是 $V = V_0 \setminus V_0 \cap H$ 上由 $F/s^{d}$ 定义的函数的情形，其中 $F$ 是次数 $d$ 的齐次多项式.有类似的素性与横截性假设，常数 $A$ 用各种 Euler-Poincar\'e 示性数表示.

这两个定理的证明占据讲义的最后两章.前三章(那些只对平展上同调有有限了解的人也能读懂)包含：关于指数和的问题(已解答与未解答的)的简短一般讨论；与指数和相联系的 $L$ 级数的回顾与可以应用 Weil 猜想的例子；解释如何用平展上同调群上的自同态的迹表示指数和，以及主要定理的陈述.

James Milne (1-MI).
